%%%%%%%%%%%%%%%%%%%%%%%%%%%%%%%%%%%%%%%%%%%%%%%%%%%%%%%%%%%%%%%%%%%%%%%%%%%%%%%%
\begin{frame}[fragile]\frametitle{}
\begin{center}
{\Large Going to Production}
\end{center}
\end{frame}
%%%%%%%%%%%%%%%%%%%%%%%%%%%%%%%%%%%%%%%%%%%%%%%%%%%%%%%%%%%%%%%%%%%%%%%%%%%%%%%%
\begin{frame}[fragile]\frametitle{Best Practices}
      \begin{itemize}
        \item \textbf{State Design}: Keep simple and clear, use type hints, only necessary information
        \item \textbf{Node Functions}: Single responsibility, handle exceptions, return only the changed fields
        \item \textbf{Edge Design}: Clear conditional logic, avoid complex cycles, consider all paths
        \item \textbf{Error Handling}: Add at critical nodes, provide fallback mechanisms, log errors
        \item \textbf{Testing}: Test individual nodes, test routing logic, test complete workflows
        \item \textbf{Observability}: Use tracing tools (Langfuse, LangSmith) for production monitoring
      \end{itemize}
\end{frame}
%%%%%%%%%%%%%%%%%%%%%%%%%%%%%%%%%%%%%%%%%%%%%%%%%%%%%%%%%%%%%%%%%%%%%%%%%%%%%%%%
\begin{frame}[fragile]\frametitle{Best Practices for Advanced Features}
      \begin{itemize}
        \item \textbf{Start Simple}: Add complexity only when needed
        \item \textbf{Checkpointing}: Always use a durable checkpointer in production
        \item \textbf{State Size}: Keep state minimal, only essential data
        \item \textbf{Cycle Limits}: Set maximum iterations to prevent infinite loops (\texttt{recursion\_limit})
        \item \textbf{Human Approval}: For high-stakes actions (financial, emails, deletions)
        \item \textbf{Documentation}: Graph visualization helps team understanding
      \end{itemize}
\end{frame}
%%%%%%%%%%%%%%%%%%%%%%%%%%%%%%%%%%%%%%%%%%%%%%%%%%%%%%%%%%%%%%%%%%%%%%%%%%%%%%%%
\begin{frame}[fragile]\frametitle{Retries and a Production Checklist}
      \begin{itemize}
        \item Attach a \texttt{RetryPolicy} to a node that calls a flaky API; LangGraph retries with exponential backoff.
      \end{itemize}
\begin{lstlisting}[language=Python]
from langgraph.types import RetryPolicy

workflow.add_node(
    "unreliable_api", unreliable_node,
    retry_policy=RetryPolicy(max_attempts=3, backoff_factor=2))
\end{lstlisting}
      \begin{itemize}
        \item \textbf{Checklist}: error handling, rate limits, tracing (LangSmith), Pydantic validation of tool inputs and outputs, per-node timeouts, strict state schema, unit tests per node.
      \end{itemize}
\end{frame}
%%%%%%%%%%%%%%%%%%%%%%%%%%%%%%%%%%%%%%%%%%%%%%%%%%%%%%%%%%%%%%%%%%%%%%%%%%%%%%%%
\begin{frame}[fragile]\frametitle{Key Features of LangGraph Studio}
      \begin{itemize}
        \item \textbf{Graph Visualization}: Your code is rendered as a graph, so you see nodes, edges and branches as they are defined.
        \item \textbf{Real-time Debugging}: Step through your graph's execution, inspect the state at each node, and identify issues quickly.
        \item \textbf{State Editing and Replay}: Edit the state at a checkpoint and re-run from that point.
        \item \textbf{Integrated LangSmith Tracing}: Seamlessly view detailed traces of your agent's runs, including LLM calls and tool usage.
        \item \textbf{Human-in-the-Loop}: Interrupts surface in the UI so you can review, approve or correct a step.
        \item \textbf{Local to Cloud}: The same project runs locally against your code, and can be deployed as an API server.
      \end{itemize}
\end{frame}
%%%%%%%%%%%%%%%%%%%%%%%%%%%%%%%%%%%%%%%%%%%%%%%%%%%%%%%%%%%%%%%%%%%%%%%%%%%%%%%%
\begin{frame}[fragile]\frametitle{Comparison: LangGraph vs. Other Agent Frameworks}
\tiny
\begin{tabular}{|l|p{3.5cm}|p{3.5cm}|p{3.5cm}|}
\hline
\textbf{Feature} & \textbf{LangGraph} & \textbf{CrewAI} & \textbf{AutoGen} \\
\hline
\textbf{Core Idea} & A library for building stateful, multi-actor applications with graphs. & A framework for orchestrating role-playing autonomous AI agents. & A framework for simplified multi-agent conversation workflows. \\
\hline
\textbf{Control} & \textbf{High}. Explicit control flow via graph structure. You define every possible path. & \textbf{Medium}. High-level process definition (e.g., sequential, hierarchical). Less granular control. & \textbf{Medium-Low}. Control is primarily through agent prompts and conversation patterns. \\
\hline
\textbf{Architecture} & Cyclic Graphs. Nodes are functions, edges define flow. Very flexible. & Agent-Task-Crew. Agents are given specific roles and tools; a process orchestrates them. & Conversation-based. Agents are "speakers" in a group chat, managed by a group chat manager. \\
\hline
\textbf{Flexibility} & \textbf{Very High}. Can build any custom flow, including human-in-the-loop and dynamic routing. & \textbf{Medium}. Optimized for common collaboration patterns but less flexible for arbitrary cycles. & \textbf{Medium}. Good for conversational agents, but complex, non-linear workflows are harder to model. \\
\hline
\textbf{Best For} & Production systems requiring reliability, control, and observability. Complex, long-running workflows. & Rapidly prototyping hierarchical or sequential agent teams for tasks like research or content creation. & Research and simulation of agent interactions and conversational systems. \\
\hline
\end{tabular}
\end{frame}
%%%%%%%%%%%%%%%%%%%%%%%%%%%%%%%%%%%%%%%%%%%%%%%%%%%%%%%%%%%%%%%%%%%%%%%%%%%%%%%%
\begin{frame}[fragile]\frametitle{Production Deployment Considerations}
      \begin{itemize}
        \item \textbf{Checkpointer Choice}:
        \begin{itemize}
            \item \texttt{MemorySaver}: For development only.
            \item \texttt{SqliteSaver}: Good for single-server, low-concurrency deployments.
            \item \texttt{PostgresSaver} / \texttt{RedisSaver}: Essential for production. They support high concurrency and data resilience, allowing you to scale your application across multiple server instances.
        \end{itemize}
        \item \textbf{Scalability}:
        \begin{itemize}
            \item Use asynchronous node functions (\texttt{async def}) and invoke the graph with \texttt{ainvoke()} or \texttt{astream()} for better performance under load.
            \item LangGraph's stateful nature allows you to run long-running tasks on background workers.
        \end{itemize}
        \item \textbf{Monitoring \& Observability}:
        \begin{itemize}
            \item \textbf{LangSmith is critical}. It provides detailed, visual traces of every graph execution, making it possible to debug complex agent behavior, track costs, and monitor performance in production.
            \item Implement logging within your nodes to capture key business events.
        \end{itemize}
      \end{itemize}
\end{frame}
%%%%%%%%%%%%%%%%%%%%%%%%%%%%%%%%%%%%%%%%%%%%%%%%%%%%%%%%%%%%%%%%%%%%%%%%%%%%%%%%
\begin{frame}[fragile]\frametitle{Serving an Agent: FastAPI and Docker}
      \begin{itemize}
        \item Wrap the compiled graph in a small web service; the client sends a \texttt{thread\_id} so each user keeps their own memory.
        \item Package it in a container for repeatable deployment.
      \end{itemize}
% [code-font: forced \scriptsize | overflow at the default size]
\begin{lstlisting}[language=Python, basicstyle=\scriptsize\ttfamily]
# server.py
from fastapi import FastAPI
from pydantic import BaseModel

api = FastAPI()

class ChatRequest(BaseModel):
    message: str
    thread_id: str = "default"

@api.post("/chat")
async def chat(req: ChatRequest):
    config = {"configurable": {"thread_id": req.thread_id}}
    result = await agent_app.ainvoke(
        {"messages": [("user", req.message)]}, config=config)
    return {"response": result["messages"][-1].content}
\end{lstlisting}
% [code-font: forced \scriptsize | overflow at the default size]
\begin{lstlisting}[basicstyle=\scriptsize\ttfamily]
FROM python:3.11-slim
WORKDIR /app
COPY requirements.txt .
RUN pip install -r requirements.txt
COPY . .
CMD ["uvicorn", "server:api", "--host", "0.0.0.0", "--port", "8000"]
\end{lstlisting}

{\tiny (Ref: Building AI Agents with LangGraph: A Practical Guide - DEV Community)}
\end{frame}
